\documentclass[10pt,a4paper]{amsart}
\usepackage[margin=19mm]{geometry}
\usepackage{amsmath,amssymb,mathtools}
\usepackage[scr=rsfs,cal=euler]{mathalfa}
\usepackage{xfrac}
\usepackage[unicode,hidelinks,bookmarks=false]{hyperref}
\newcommand{\ie}{i.e.\ }
\newcommand{\cf}{cf.\ }
\newcommand{\resp}{respectively\ }
\newcommand{\Subsection}{\subsection}
\providecommand{\nobreakdash}{\nobreak\hbox{-}}
\setlength{\emergencystretch}{2em}
\multlinegap0pt
\usepackage{amssymb}
\usepackage{bm}
\usepackage{tensor}
\usepackage[shortlabels]{enumitem}
\usepackage{xcolor}
\newtheorem{proposition}{Proposition}
\newcommand{\wed}{\wedge\hspace*{-7.5pt}.\hspace*{3pt}}
\title{Hidden diffeos in the Hamiltonian formulation of a background independent field theory}
\author{J. Fernando Barbero G., Bogar D\'{\i}az, Juan Margalef-Bentabol, Eduardo J.S. Villase\~nor}
\date{Source: arXiv:2507.12184v1 (CC BY 4.0)}
\begin{document}
\maketitle

\noindent\emph{Source and licence.} Hidden diffeos in the Hamiltonian formulation of a background independent field theory. Version arXiv:2507.12184v1 (CC BY 4.0), distributed by arXiv under the Creative Commons Attribution 4.0 International licence, http://creativecommons.org/licenses/by/4.0/, as recorded in the arXiv licence metadata for this exact version and shown on the abstract page https://arxiv.org/abs/2507.12184v1. Full licence notice: \texttt{licenses/arxiv-2507.12184-CC-BY-4.0.txt}.\par\smallskip

\noindent\textit{Editorial scope.} Counted unit: the article's proposition thm\_Lagrangian (source0.tex lines 450--456). The article's complete original proof of it is delivered with it (source0.tex lines 457--478, one \texttt{proof} environment of 22 lines). No mathematics is written, repaired or re-derived: the statement and the proof are the article's own lines, in the article's own notation, and no retained line is substituted or re-flowed.  Retained uncounted as the source-local context the proof needs: lines 350--352; lines 443--448.  Nothing was omitted from any retained span, so the package carries no omission record.  No retained line contains a citation, so the wrapper carries no bibliography and no bibliography entry.  No retained line contains a figure environment, an image-inclusion command or drawing code, so no image is delivered and the package declares no asset dependency.  Editorial formatting only, in addition: an AMS \texttt{amsart} preamble that restates the article's own declarations and macros used by the retained lines, namely the environments \texttt{proposition} on separate counters, together with the standard packages those lines need.  The retained environments print in this document as Proposition 1 in order of appearance, whereas the article prints the same items with the numbering of its own full document.  The complete native source of the article as distributed in the arXiv e-print for this exact version is bundled unchanged as \texttt{sources/181679/source0.tex}.
\begin{equation}\label{action_HM}
S_{\mathrm{HM}}(\Phi,\mathrm{A})=\frac{1}{2}\int_{[\tau_1,\tau_2]\times \Sigma}\langle[\mathrm{d}_{\mathrm{A}}\Phi\wed \mathrm{d}_{\mathrm{A}}\Phi]\wed \mathrm{F}_{\mathrm{A}}\rangle\,.
\end{equation}

\begin{align}\label{Lagrangian1}
\begin{split}
  \mathsf{L}(v_q)&=\int_\Sigma \Big(\langle [v_\phi\wed \mathrm{d}_A\phi]\wed F_A\rangle +\langle [[a\wed \phi]\wed \mathrm{d}_A\phi]\wed F_A\rangle \\
   &\hspace*{42mm}+\frac{1}{2}\langle [\mathrm{d}_A\phi \wed \mathrm{d}_A\phi]\wed v_A\rangle-\frac{1}{2}\langle  [\mathrm{d}_A\phi \wed \mathrm{d}_A\phi]\wed \mathrm{d}_A a \rangle\Big)
\end{split}
\end{align}

\begin{proposition}
    \label{thm_Lagrangian}
The Lagrangian \eqref{Lagrangian1} defined by the action \eqref{action_HM} is
\begin{equation}\label{Lagrangian2}
 \mathsf{L}(v_q)=\int_\Sigma \Big(\langle [v_\phi\wed \mathrm{d}_A\phi]\wed F_A\rangle +\frac{1}{2}\langle v_A\wed[\mathrm{d}_A\phi \wed \mathrm{d}_A\phi]\rangle+\langle a\wed \mathrm{d}_A[[F_A\wed\phi]\wed\phi]\rangle \Big)\,.
\end{equation}
\end{proposition}

\begin{proof}
Indeed, we have
\begin{align*}
\langle [[ a\wed \phi]\wed \mathrm{d}_A\phi]\wed F_A \rangle&=\langle [ a\wed\phi ] \wed [ \mathrm{d}_A\phi\wed F_A]\rangle = \langle  a\wed[\phi \wed [\mathrm{d}_A \phi\wed F_A]]\rangle\\
&=-\langle a\wed[[\mathrm{d}_A\phi\wed F_A]\wed\phi] \rangle=\langle a\wed[[F_A\wed \mathrm{d}_A\phi]\wed\phi] \rangle
\end{align*}
Also
\begin{align*}
0=\int_\Sigma -\frac{1}{2}\mathrm{d}_A\langle [\mathrm{d}_A\phi\wed \mathrm{d}_A\phi]\wed a   \rangle&=\int_\Sigma \Big( -\langle[[F_A\wed \phi]\wed \mathrm{d}_A\phi ]\wed a \rangle-\frac{1}{2}\langle[\mathrm{d}_A\phi\wed \mathrm{d}_A\phi]\wed \mathrm{d}_A a\rangle  \Big)
\end{align*}
so that
\[
\int_\Sigma -\frac{1}{2}\langle[\mathrm{d}_A\phi\wed \mathrm{d}_A\phi]\wed \mathrm{d}_A a\rangle=\int_\Sigma\langle a\wed[[F_A\wed \phi]\wed \mathrm{d}_A\phi]\rangle\,.
\]
Finally, we conclude that the two velocity-independent terms in \eqref{Lagrangian1} add up to
\begin{align*}
&\int_\Sigma \Big(\langle [[a\wed \phi]\wed \mathrm{d}_A\phi]\wed F_A\rangle-\frac{1}{2}\langle  [\mathrm{d}_A\phi \wed \mathrm{d}_A\phi]\wed \mathrm{d}_A a \rangle\Big)\\
=&\int_\Sigma\Big(\langle a\wed[[F_A\wed \mathrm{d}_A\phi]\wed \phi]\rangle+\langle a\wed[[F_A\wed\phi]\wed \mathrm{d}_A\phi]  \rangle\Big)\\
=&\int_\Sigma\Big(\langle a\wed[\mathrm{d}_A[F_A\wed\phi]\wed \phi]\rangle+\langle a\wed[[F_A\wed\phi]\wed \mathrm{d}_A\phi]  \rangle\Big)\\
=&\int_\Sigma \langle a\wed \mathrm{d}_A[[F_A\wed \phi]\wed \phi]\rangle\,.
\end{align*}    
\end{proof}

\end{document}
